% =====================================================================
% Skeleton G — 3 Layer Vertical + Central Hero + N Parallel Instances
% =====================================================================
% USAGE: Adapted from a GOLD STANDARD figure (Federated Traffic Prediction).
% Sub-agent: copy ENTIRE file as figure.tex, change content.
%
% ⚙ LAYOUT = BY-CONSTRUCTION (positioning + fit + node anchors): the 3
%   city ST-GNN zones are the spine (city pitch = one macro), the server
%   sits at a calc midpoint above them, prep boxes are below=of their city,
%   sensors offset from prep, upload/download arrows ride node anchors.
%   Change content / add a city and the rest REFLOWS — no coordinate math.
%   ⚠ The "zone bottom y=8.6 / 4.6 / chart y=0.4–4.6" numbers below are
%   ORIGINAL design intent — zone bounds and chip clearance are now
%   maintained BY CONSTRUCTION (chips anchored above each zone .north;
%   chart panel via fit). SEMANTIC constraints (renaming, small_blue ≤5-CJK
%   text limit, recompute-bar-heights-with-labels) STILL HOLD.
%
% LAYOUT MODEL:
%   ┌────────────────────────────────────────────────┐  ┌───┐
%   │   Top Layer: Central Hero (aggregator/server)  │  │ R │
%   │              [Big purple box]                  │  │ i │
%   ├────────────────────────────────────────────────┤  │ g │
%   │ Mid Layer: N parallel instances (3 cities/    │  │ h │
%   │  workers/clients) each with internal hero      │  │ t │
%   ├────────────────────────────────────────────────┤  │   │
%   │ Bot Layer: Edge/Input + multi-modal sources    │  │ V │
%   │ (camera/sensor/GPS for each instance)          │  │ . │
%   └────────────────────────────────────────────────┘  │ l │
%   Bottom-most: experiment result bar chart            │ . │
%                                                       └───┘
%
% WHEN TO USE: Federated learning, distributed multi-agent systems,
% edge-cloud architectures. Use when you have 1 central + N parallel.
%
% CONTAINS: Central hero (FedAvg server) + 3 parallel cities each with
% ST-GNN (spatial attn + temporal conv) + edge preprocessing per city +
% multi-modal sensors (camera/磁/GPS) + result bar chart + dashed
% feedback rails (上传/下发) + 中文 vertical right label "联邦迭代".
%
% CONSTRAINTS — DO NOT VIOLATE (each one fixed a real bug):
%
%   [RENAMING]
%   - 中心服务器 / 加权 FedAvg 聚合 — RENAME for non-FL topics
%     (e.g., "Coordinator / Multi-Agent Aggregator" for multi-agent LLM).
%   - 3 city names + 3 ST-GNN heroes — RENAME (e.g., 3 agents/workers).
%   - 3 sets of multi-modal sensors (camera/磁/GPS) — RENAME to topic
%     input modalities.
%   - "联邦迭代" right-side vertical label — RENAME to topic loop name.
%   - Bottom bar chart labels AND bar heights — RENAME labels AND
%     recompute bar heights to match the new numeric labels.
%
%   [small_blue BOX TEXT LIMIT — hardcoded 1.6cm width]
%   - small_blue is fixed 1.6cm wide → max ~5 CJK chars OR ~10 English
%     chars. Text like "LLM (GPT-4)" truncates inside the box.
%     Use a single short token ("LLM" or "GPT-4"), not both.
%
%   [ZONE CHIPS (3 layer labels) — position above zones]
%   - All 3 chips use anchor=south west, positioned ABOVE their zone
%     (chip y > zone top y). Do NOT place chips inside the zone's
%     top-left corner — they overlap the content immediately below.
%
%   [CHART PANEL — must enclose ALL its parts]
%   - Panel rectangle MUST span the full y-range of: x-axis arrow tip,
%     AND axis tick labels, AND chart title. Current bounds: y=0.4 to
%     y=4.6 (height 4.2cm); title at y=3.9; "数据不出域" chip at y=3.9.
%   - Bar heights MUST match their numeric labels — recompute if you
%     change data values, do NOT keep old heights.
%
%   [ZONE BOUNDS — leave room for adjacent chips/titles]
%   - Training zone bottom y=8.6 (gap above for 本地训练层 chip).
%   - Edge zone bottom y=4.6 (clearance for chart title at y=3.9).
% =====================================================================

\documentclass[tikz,border=30pt]{standalone}
\usepackage{tikz}
\usepackage[fontset=none]{ctex}
\setCJKmainfont{PingFang SC}
\setCJKsansfont{PingFang SC}
\usepackage{amsmath}
\usetikzlibrary{shapes, arrows.meta, positioning, fit, backgrounds, calc, shadows}

% ===== Color palette =====
\definecolor{drawBlueFill}{HTML}{DAE8FC}    \definecolor{drawBlueLine}{HTML}{6C8EBF}
\definecolor{drawGreenFill}{HTML}{D5E8D4}   \definecolor{drawGreenLine}{HTML}{82B366}
\definecolor{drawOrangeFill}{HTML}{FFE6CC}   \definecolor{drawOrangeLine}{HTML}{D79B00}
\definecolor{drawPurpleFill}{HTML}{E1D5E7}   \definecolor{drawPurpleLine}{HTML}{9673A6}
\definecolor{drawRedFill}{HTML}{F8CECC}      \definecolor{drawRedLine}{HTML}{B85450}
\definecolor{drawGreyFill}{HTML}{F5F5F5}     \definecolor{drawGreyLine}{HTML}{666666}
\definecolor{barBlue}{HTML}{3B82F6}

\pgfdeclarelayer{background}
\pgfsetlayers{background,main}

% ---- single source of truth for layout rhythm ----
\def\citegap{1.6cm}   % horizontal gap BETWEEN the 4.4cm city zones (=> 6cm pitch)

\begin{document}
\begin{tikzpicture}[
    every node/.style={font=\small},
    base_box/.style={rectangle, rounded corners=3pt, align=center,
        minimum height=0.9cm, minimum width=2.8cm,
        inner sep=10pt, drop shadow={opacity=0.15}, thick},
    green_node/.style={base_box, fill=drawGreenFill, draw=drawGreenLine},
    red_node/.style={base_box, fill=drawRedFill, draw=drawRedLine, font=\small\bfseries},
    grey_node/.style={base_box, fill=drawGreyFill, draw=drawGreyLine},
    small_blue/.style={rectangle, rounded corners=2pt, align=center,
        minimum height=0.65cm, minimum width=1.6cm,
        inner sep=5pt, thick, font=\footnotesize,
        fill=drawBlueFill!60, draw=drawBlueLine!80},
    sensor/.style={rectangle, rounded corners=2pt, align=center,
        minimum height=0.5cm, minimum width=1.2cm,
        inner sep=3pt, thick, font=\scriptsize,
        fill=drawOrangeFill, draw=drawOrangeLine},
    arrow/.style={-{Stealth[scale=0.9]}, thick, color=black!70,
        shorten >=2pt, shorten <=1pt},
    thin_arrow/.style={-{Stealth[scale=0.7]}, thick, color=drawGreyLine},
    zone label/.style={font=\footnotesize\bfseries, rounded corners=3pt,
        inner sep=4pt, anchor=south west},
]

% ================================================================
%  LAYER 2 — LOCAL TRAINING (middle) — the spine of the layout.
%  Place the 3 city zones first; everything else hangs off them.
%  City A anchored; B, C chained `right=of` previous (pitch = \citegap).
% ================================================================
\node[green_node, minimum width=4.4cm, minimum height=3.2cm] (localA) at (0, 10.5) {};
\node[green_node, minimum width=4.4cm, minimum height=3.2cm,
      right=\citegap of localA] (localB) {};
\node[green_node, minimum width=4.4cm, minimum height=3.2cm,
      right=\citegap of localB] (localC) {};

% City internal blocks — each anchored to its own city node so it travels
% with the column.  Titles near the top, spat/temp pair mid, caption low.
\node[font=\small\bfseries, text=drawGreenLine!80!black]
    at ($(localA.north)+(0,-1.1)$) {城市A\,$\cdot$\,ST-GNN};
\node[small_blue] (spatA) at ($(localA.center)+(-0.8,-0.3)$) {空间注意力};
\node[small_blue] (tempA) at ($(localA.center)+(0.8,-0.3)$) {时间卷积};
\draw[thin_arrow] (spatA.east) -- (tempA.west);
\node[font=\scriptsize, text=drawGreyLine, align=center]
    at ($(localA.south)+(0,0.7)$) {路网拓扑\,$\to$\,周期模式};

\node[font=\small\bfseries, text=drawGreenLine!80!black]
    at ($(localB.north)+(0,-1.1)$) {城市B\,$\cdot$\,ST-GNN};
\node[small_blue] (spatB) at ($(localB.center)+(-0.8,-0.3)$) {空间注意力};
\node[small_blue] (tempB) at ($(localB.center)+(0.8,-0.3)$) {时间卷积};
\draw[thin_arrow] (spatB.east) -- (tempB.west);
\node[font=\scriptsize, text=drawGreyLine, align=center]
    at ($(localB.south)+(0,0.7)$) {路网拓扑\,$\to$\,周期模式};

\node[font=\small\bfseries, text=drawGreenLine!80!black]
    at ($(localC.north)+(0,-1.1)$) {城市C\,$\cdot$\,ST-GNN};
\node[small_blue] (spatC) at ($(localC.center)+(-0.8,-0.3)$) {空间注意力};
\node[small_blue] (tempC) at ($(localC.center)+(0.8,-0.3)$) {时间卷积};
\draw[thin_arrow] (spatC.east) -- (tempC.west);
\node[font=\scriptsize, text=drawGreyLine, align=center]
    at ($(localC.south)+(0,0.7)$) {路网拓扑\,$\to$\,周期模式};

% ================================================================
%  LAYER 3 — FEDERATED AGGREGATION (top)
%  Server centered over the city row via calc midpoint, fixed height above.
% ================================================================
\coordinate (cityrowtop) at ($(localA.north)!0.5!(localC.north)$);
\node[red_node, minimum width=7.5cm, minimum height=2.4cm] (server)
    at ($(cityrowtop)+(0,4.4)$)
    {\textbf{中心服务器}\\[2pt]{\footnotesize 加权 FedAvg 聚合}\\[-1pt]{\footnotesize 生成并下发全局模型}};

% ================================================================
%  LAYER 1 — EDGE COLLECTION (bottom)
%  Preprocessing box below each city; sensors below each prep box.
% ================================================================
\node[grey_node, minimum width=3.2cm, below=1.45cm of localA] (prepA)
    {\footnotesize 本地预处理};
\node[grey_node, minimum width=3.2cm, below=1.45cm of localB] (prepB)
    {\footnotesize 本地预处理};
\node[grey_node, minimum width=3.2cm, below=1.45cm of localC] (prepC)
    {\footnotesize 本地预处理};

% Sensors — 3 per city, fanned under each prep box (offset -1.5/0/+1.5)
\node[sensor] (camA) at ($(prepA.south)+(-1.5,-1.45)$) {摄像头};
\node[sensor] (magA) at ($(prepA.south)+(0,-1.45)$) {地磁};
\node[sensor] (gpsA) at ($(prepA.south)+(1.5,-1.45)$) {GPS};
\node[sensor] (camB) at ($(prepB.south)+(-1.5,-1.45)$) {摄像头};
\node[sensor] (magB) at ($(prepB.south)+(0,-1.45)$) {地磁};
\node[sensor] (gpsB) at ($(prepB.south)+(1.5,-1.45)$) {GPS};
\node[sensor] (camC) at ($(prepC.south)+(-1.5,-1.45)$) {摄像头};
\node[sensor] (magC) at ($(prepC.south)+(0,-1.45)$) {地磁};
\node[sensor] (gpsC) at ($(prepC.south)+(1.5,-1.45)$) {GPS};

% ================================================================
%  EXPERIMENT RESULTS PANEL (bottom)
%  Panel centered under the city row; all internals kept verbatim inside
%  a scope shifted to the panel's lower-left so the embedded chart and
%  legend travel with the panel.
% ================================================================
\node[rectangle, rounded corners=6pt, draw=drawBlueLine, fill=drawBlueFill!15,
    minimum width=15cm, minimum height=4.2cm, thick, inner sep=0pt] (panel)
    at ($(cityrowtop |- 0,1.7)$) {};

% Title + key-insight chip anchored to the panel's TOP edge so they stay
% aligned with it (above the panel, anchor=south).
\node[font=\small\bfseries, text=drawBlueLine!80!black, anchor=south west]
    at ($(panel.north west)+(0.0,0.1)$) {实验结果：三城市联合预测性能};
\node[font=\footnotesize\bfseries, text=drawBlueLine!80!black,
      fill=drawBlueFill!40, draw=drawBlueLine!50, rounded corners=3pt,
      inner sep=5pt, anchor=south east] at ($(panel.north east)+(-0.2,0.1)$)
    {数据不出域，隐私安全};

% Embedded bar chart + legend — original LOCAL coordinates preserved,
% just re-anchored: the original used absolute origin near panel SW corner
% (panel center (6,1.7) => SW corner (-1.5,-0.4)). Shift the scope to the
% panel SW corner and keep the exact internal drawing.
\begin{scope}[shift={($(panel.south west)+(1.5,0.0)$)}]
  % --- bar chart (was shift={(1.5,0.20)} relative to abs origin; here the
  %     scope already sits at panel SW + (1.5,0), so use +0.20 in y) ---
  \begin{scope}[shift={(0, 0.20)}]
    % Axes — extended to y=3.2 (label 30)
    \draw[-{Stealth[scale=0.6]}, thick, drawGreyLine] (0.5, 0.2) -- (0.5, 3.4);
    \draw[-{Stealth[scale=0.6]}, thick, drawGreyLine] (0.5, 0.2) -- (5.5, 0.2);
    % Grid lines + Y-axis ticks at 5/10/15/20/25/30
    \foreach \y/\lab in {0.7/5, 1.2/10, 1.7/15, 2.2/20, 2.7/25, 3.2/30} {
      \draw[gray!20, thin] (0.5, \y) -- (5.2, \y);
      \node[font=\scriptsize, left, text=drawGreyLine] at (0.4, \y) {\lab};
    }
    % MAE group — bar heights match axis (22.1->y=2.41, 18.1->y=2.01)
    \fill[drawGreyFill, draw=drawGreyLine] (1.0, 0.2) rectangle (1.6, 2.41);
    \fill[barBlue, draw=barBlue!80!black, rounded corners=1pt] (1.8, 0.2) rectangle (2.4, 2.01);
    \node[font=\scriptsize, text=drawGreyLine] at (1.3, 2.56) {22.1};
    \node[font=\scriptsize, text=barBlue!80!black] at (2.1, 2.16) {18.1};
    \node[font=\scriptsize] at (1.7, -0.05) {MAE};
    % RMSE group — 25.3->y=2.73, 19.8->y=2.18
    \fill[drawGreyFill, draw=drawGreyLine] (3.0, 0.2) rectangle (3.6, 2.73);
    \fill[barBlue, draw=barBlue!80!black, rounded corners=1pt] (3.8, 0.2) rectangle (4.4, 2.18);
    \node[font=\scriptsize, text=drawGreyLine] at (3.3, 2.88) {25.3};
    \node[font=\scriptsize, text=barBlue!80!black] at (4.1, 2.33) {19.8};
    \node[font=\scriptsize] at (3.7, -0.05) {RMSE};
    % Improvement arrows + labels — to the RIGHT of each pair
    \draw[-{Stealth[scale=0.5]}, very thick, drawRedLine] (2.55, 2.41) -- (2.55, 2.01);
    \node[font=\scriptsize\bfseries, text=drawRedLine, anchor=west]
        at (2.6, 2.21) {$\downarrow$18.3\%};
    \draw[-{Stealth[scale=0.5]}, very thick, drawRedLine] (4.55, 2.73) -- (4.55, 2.18);
    \node[font=\scriptsize\bfseries, text=drawRedLine, anchor=west]
        at (4.6, 2.46) {$\downarrow$21.7\%};
  \end{scope}
  % --- legend (was shift={(8.5,1.5)} relative to abs origin; here scope is
  %     at panel SW + (1.5,0) so legend origin = (7.0, 1.5)) ---
  \begin{scope}[shift={(7.0, 1.5)}]
    \fill[drawGreyFill, draw=drawGreyLine] (0, 0.7) rectangle (0.4, 0.9);
    \node[font=\scriptsize, anchor=west] at (0.5, 0.8) {基线方法};
    \fill[barBlue, draw=barBlue!80!black] (0, 0.25) rectangle (0.4, 0.45);
    \node[font=\scriptsize, anchor=west] at (0.5, 0.35) {FedTraffic};
  \end{scope}
\end{scope}

% ================================================================
%  CONNECTIONS — Sensors -> Preprocessing (fan-in, NODE ANCHORS)
% ================================================================
\draw[arrow] (camA.north) -- ([xshift=-10pt]prepA.south);
\draw[arrow] (magA.north) -- (prepA.south);
\draw[arrow] (gpsA.north) -- ([xshift=10pt]prepA.south);
\draw[arrow] (camB.north) -- ([xshift=-10pt]prepB.south);
\draw[arrow] (magB.north) -- (prepB.south);
\draw[arrow] (gpsB.north) -- ([xshift=10pt]prepB.south);
\draw[arrow] (camC.north) -- ([xshift=-10pt]prepC.south);
\draw[arrow] (magC.north) -- (prepC.south);
\draw[arrow] (gpsC.north) -- ([xshift=10pt]prepC.south);

% ================================================================
%  CONNECTIONS — Preprocessing -> Local Models (NODE ANCHORS)
% ================================================================
\draw[arrow] (prepA.north) -- (localA.south);
\draw[arrow] (prepB.north) -- (localB.south);
\draw[arrow] (prepC.north) -- (localC.south);

% ================================================================
%  UPLOAD GRADIENTS (orange thick, L-shaped convergence)
%  Same L-shape as original; corners land on the server's SOUTH edge
%  (node anchor + xshift) instead of re-typed (4.0,15.3)/(8.0,15.3).
% ================================================================
% A -> server: up, right, up  (target = server.south shifted left 2.0cm)
\draw[-{Stealth[scale=1.1]}, very thick, color=drawOrangeLine, rounded corners=6pt]
    ([xshift=-10pt]localA.north) -- ++(0, 1.5) -| ([xshift=-57pt]server.south);
% B -> server: straight up
\draw[-{Stealth[scale=1.1]}, very thick, color=drawOrangeLine]
    ([xshift=-10pt]localB.north) -- ([xshift=-10pt]server.south);
% C -> server: up, left, up  (target = server.south shifted right 2.0cm)
\draw[-{Stealth[scale=1.1]}, very thick, color=drawOrangeLine, rounded corners=6pt]
    ([xshift=-10pt]localC.north) -- ++(0, 1.5) -| ([xshift=57pt]server.south);

% Upload label
\node[font=\footnotesize\bfseries, text=drawOrangeLine,
      fill=white, fill opacity=0.9, text opacity=1,
      inner sep=3pt, rounded corners=2pt]
    at ($(localA.north)!0.5!(localB.north)+(0,1.6)$) {上传梯度};

% ================================================================
%  DOWNLOAD MODEL (blue dashed, L-shaped divergence)
% ================================================================
% server -> A: down, left, down  (start = server.south shifted left 1.5cm)
\draw[-{Stealth[scale=0.9]}, thick, dashed, color=drawBlueLine, rounded corners=6pt]
    ([xshift=-43pt]server.south) -- ++(0, -1.3) -| ([xshift=10pt]localA.north);
% server -> B: straight down
\draw[-{Stealth[scale=0.9]}, thick, dashed, color=drawBlueLine]
    ([xshift=10pt]server.south) -- ([xshift=10pt]localB.north);
% server -> C: down, right, down  (start = server.south shifted right 1.5cm)
\draw[-{Stealth[scale=0.9]}, thick, dashed, color=drawBlueLine, rounded corners=6pt]
    ([xshift=43pt]server.south) -- ++(0, -1.3) -| ([xshift=10pt]localC.north);

% Download label
\node[font=\footnotesize\bfseries, text=drawBlueLine,
      fill=white, fill opacity=0.9, text opacity=1,
      inner sep=3pt, rounded corners=2pt]
    at ($(localB.north)!0.5!(localC.north)+(0,1.6)$) {下发全局模型};

% ================================================================
%  ZONE BACKGROUNDS — drawn via `fit` over the real member nodes so the
%  zone box follows them under reflow.  inner sep tuned to reproduce the
%  original dashed-rectangle margins.
% ================================================================
\begin{pgfonlayer}{background}
  % Aggregation zone — wraps the server; wide margin matches the original
  % full-width band (orig spanned the city-row width, ~1cm beyond cities).
  \node[fit=(server), inner xsep=155pt, inner ysep=18pt,
        fill=drawRedFill!20, rounded corners=10pt,
        draw=drawRedLine!50, dashed, thick] (aggzone) {};
  % Training zone — wraps all 3 city zones (~1cm horizontal margin).
  \node[fit=(localA)(localC), inner xsep=28pt, inner ysep=12pt,
        fill=drawGreenFill!20, rounded corners=10pt,
        draw=drawGreenLine!50, dashed, thick] (trainzone) {};
  % Edge collection zone — wraps prep boxes + all sensors (~1cm margin).
  \node[fit=(prepA)(prepC)(camA)(gpsC), inner xsep=28pt, inner ysep=12pt,
        fill=drawOrangeFill!20, rounded corners=10pt,
        draw=drawOrangeLine!50, dashed, thick] (edgezone) {};
\end{pgfonlayer}

% ================================================================
%  ZONE LABELS (capsule chips — placed ABOVE each zone via anchors)
% ================================================================
\node[zone label, text=drawRedLine!80!black,
      fill=drawRedFill!60, draw=drawRedLine!40]
    at ($(aggzone.north west)+(0.0,0.05)$) {联邦聚合层};
\node[zone label, text=drawGreenLine!80!black,
      fill=drawGreenFill!60, draw=drawGreenLine!40]
    at ($(trainzone.north west)+(0.0,0.05)$) {本地训练层};
\node[zone label, text=drawOrangeLine!80!black,
      fill=drawOrangeFill!60, draw=drawOrangeLine!40, inner sep=3pt]
    at ($(edgezone.north west)+(0.0,0.05)$) {边缘采集层};

% ================================================================
%  DECORATIVE — Iteration cycle annotation (right margin)
%  Anchored to the aggregation/training zones on the right edge.
% ================================================================
\coordinate (railtop) at ($(aggzone.east)+(0.1,0)$);
\coordinate (railbot) at ($(trainzone.east)+(0.1,0)$);
\draw[-{Stealth[scale=0.8]}, thick, color=drawPurpleLine, rounded corners=4pt]
    (railtop) -- (railbot);
\draw[-{Stealth[scale=0.8]}, thick, dashed, color=drawPurpleLine, rounded corners=4pt]
    ($(railbot)+(0.4,0)$) -- ($(railtop)+(0.4,0)$);
\node[font=\scriptsize, text=drawPurpleLine, anchor=west]
    at ($(railtop)!0.5!(railbot)+(0.7,0)$)
    {\shortstack{联\\邦\\迭\\代}};

\end{tikzpicture}
\end{document}
